Orbiter can work with the Grassmannian over a finite field. 
In particular, Orbiter offers indexing for subspaces.
For the special case of the Grassmannian ${\mathfrak Gr}_{4,2}(V)$, 
Pl\"ucker coordinates can be used to identify 
${\mathfrak Gr}_{4,2}(V)$ with the $Q^+(5,q)$ (Klein) quadric.
Here is an example. The command 

\sourcecode{GR_4_2_2}

creates the elements of ${\mathfrak Gr}_{4,2,2}$ 
and lists them together with their Pl\"ucker coordinates.
The following list is produced (output shortened): 
\orbiteroutput{GRAPHICS/LATEX/report_Gr_4_2_2}
The Pl\"ucker coordinates satisfy
$$
p_{12}p_{34} - p_{13}p_{24}+p_{14}p_{23}=0
$$
and hence belong to the Klein quadric $Q^+(5,q)$.
Orthogonal spaces and quadrics will 
be discussed in Section~\ref{sec:orthogonal}.

\bigskip


The Orbiter labeling of points of 
the $Q^+(5,q)$ quadric (see Section~\ref{sec:orthogonal}) 
can then be used to enumerate the lines of $\PG(3,q)$ in a second, different way.
In the example of $\PG(3,2),$ 
this yields the following list (output shortened):
\orbiteroutput{GRAPHICS/LATEX/report_Op_5_2}


